\section{بحث}
\label{sec:disc}
\begin{itemize}
\item \textbf{جداسازی:} متغیرهای زیرساختار که ساختار دو‌هسته‌ای را می‌سنجند (جرم، $m_{sd}$، \lr{girth}، $\tau_{21}$، $C_2$، $D_2$، $N_2$، $R_g$) \lr{Z$\to$b$\bar{\mathrm b}$} را از \bb جدا می‌کنند. $\tau_{32}$ و $z_g$ و \pt تقریباً بی‌اثرند.
\item \textbf{انتقال‌پذیری:} در متغیرهای جداکننده، نسبت سیگنال به پس‌زمینه برای دو جفت تقریباً یکی است. بنابراین شناساگری که روی یک جفت آموزش دیده در این متغیرها روی جفت دیگر هم باید کار کند. ولی توزیع همین متغیرها برای خود سیگنال در دو فرایند یکی نیست: \eta و قطبش \lr{Z} فرق دارند (بخش~\ref{sec:zorigin})، پس این انتظار باید با آزمون واقعی انتقال بررسی شود.
\item \textbf{خطر میان‌بر:} $\eta$ در دو جفت رفتار برعکس دارد (بخش~\ref{sec:set4}). متغیرهای کینماتیکی مانند $\eta$ و $p_z$ و $E$ و $\cos\theta$ ممکن است به‌جای ماهیت \lr{Z}، نمونه را شناسایی کنند.
\item \textbf{محدودیت پس‌زمینه:} هر دو پس‌زمینه فرایندهای \lr{QCD} خالص‌اند (\lr{bbbb} با \lr{QED=0} و \ttbb با \lr{WEIGHTED<=4}، یعنی بدون دیاگرام \lr{Z}، \lr{$\gamma^*$} و هیگز)، پس زوج \bb حاصل برهمکنش \lr{QCD} است. اینکه زوج از شکافت گلوئون می‌آید یا مستقیم از فرایند سخت را جدا نکرده‌ایم. جت سبک و جت تاپ در پس‌زمینه نیامده است.
\item \textbf{برش‌های پارتونی نمونه‌ها یکسان نیست} (بخش~\ref{sec:sim}). پس \pt و $\eta$ باید وزن‌دهی مجدد شوند یا از ویژگی نسبی استفاده شود.
\item \textbf{بازتولیدپذیری:} در نسخه نهایی داده، یک بذر تصادفی یکسان در همه مراحل شبیه‌سازی گذاشته شده است و دو اجرای یکسان از یک کار آزمایشی (۱۰۰ رویداد) در همه متغیرهای بررسی‌شده دقیقاً یکی بودند. در مقایسه با اجرای قبلی همان کار (که بذر \lr{Delphes} ثابت نبود) تعداد جت‌ها حدود ۰.۵\% فرق دارد، که با نوسان آماری سازگار است.
\item \textbf{تعریف پس‌زمینه جزئی} (\lr{nb\_notop $\ge1$}) هنوز با استاد تأیید نشده است.
\end{itemize}

\section{گام بعدی}
\label{sec:next}
\begin{enumerate}
\item خط پایه با برش روی جرم، $\tau_{21}$ و $D_2$؛
\item استفاده از ویژگی‌های نسبی و وزن‌دهی مجدد \pt و $\eta$؛
\item آزمون \lr{leave-one-pair-out}: آموزش روی یک جفت و ارزیابی روی جفت دیگر؛
\item آموزش مشترک (\lr{joint}) روی چهار نمونه، و در صورت لزوم آموزش \lr{domain-adversarial}؛
\item جداسازی \lr{train/test} بر اساس \lr{event\_id}.
\end{enumerate}
